\chapter{Wavelet Analysis}

\section{Table of Contents}

\begin{itemize}
  \item Core Idea
  \item Key Concepts
  \item Theory and Formulas
  \item Visual Schemes
  \item Examples
\end{itemize}

\section{Core Idea}

\textbf{Wavelet analysis} represents signals with basis functions that change scale and position. Unlike DFT or block-DCT, wavelets use \textbf{adaptive multiresolution}:

\begin{itemize}
  \item low frequencies use long windows and good frequency resolution,
  \item high frequencies use short windows and good spatial/time resolution,
  \item smooth trends are separated from sharp anomalies such as edges.
\end{itemize}

For image compression, this gives sparse coefficients without $8x8$ block boundaries. JPEG2000 uses wavelet transforms, bitplane coding, arithmetic coding, and optimized truncation to provide quality scalability, resolution scalability, ROI coding, exact bitrate control, and lossy-to-lossless coding.

\begin{calloutbox}{red!70!black}{Main wavelet compression idea}
Images are modeled as smooth trends plus localized anomalies. Wavelets represent trends in low-pass subbands and edges in sparse high-pass subbands, making coefficient coding efficient.
\end{calloutbox}

\section{Key Concepts}

\subsection{Linear Transform by Projection}

All transforms used in image coding project the signal onto basis functions:

\[
x(t)=\sum_k c_k\phi_k(t)
\]

with:

\[
c_k=\langle x(t),\phi_k(t)\rangle
=
\int_{-\infty}^{+\infty}x(t)\phi_k^*(t)\,dt
\]

Large $|c_k|$ means the signal strongly resembles basis function $k$. Compression works when most coefficients are small.

\subsection{Time-Frequency Uncertainty}

For a windowed sinusoid:

\[
x[n]=\cos(\omega_0 n)\operatorname{rect}_N[n]
\]

window length controls localization:

\begin{center}\fitwidth{%
\begin{tabular}{|l|l|l|}
\hline
\textbf{Window} & \textbf{Time localization} & \textbf{Frequency resolution} \\ \hline
short & excellent & poor \\ \hline
medium & good & medium \\ \hline
long & poor & good \\ \hline
very long & very poor & excellent \\ \hline
\end{tabular}}
\end{center}

\begin{calloutbox}{red!70!black}{Time-frequency uncertainty}
\[
\Delta t\cdot\Delta f\geq\frac{1}{4\pi}
\]

The area of the time-frequency uncertainty cell is fixed. A transform can change the cell shape, but cannot make both time and frequency resolution arbitrarily good.
\end{calloutbox}

\subsection{STFT vs. Wavelets}

\textbf{STFT / block DCT} use a fixed window for all frequencies. This creates rigid tiling of the time-frequency plane.

\textbf{Wavelets} adapt the window:

\begin{center}\fitwidth{%
\begin{tabular}{|l|l|l|}
\hline
\textbf{Frequency} & \textbf{Wavelet window} & \textbf{Captures} \\ \hline
high & short & edges, impulses, anomalies \\ \hline
low & long & trends, smooth texture \\ \hline
\end{tabular}}
\end{center}

Wavelet basis from a mother wavelet:

\[
\psi_{a,b}(t)
=
\frac{1}{\sqrt{a}}\psi\left(\frac{t-b}{a}\right)
\]

where $a$ controls scale and $b$ controls position.

\begin{itemize}
  \item Large $a$: stretched wavelet, lower frequency, better frequency resolution.
  \item Small $a$: compressed wavelet, higher frequency, better time/spatial resolution.
\end{itemize}

\subsection{Image Model}

Images contain:

\begin{center}\fitwidth{%
\begin{tabular}{|l|l|l|l|}
\hline
\textbf{Component} & \textbf{Signal behavior} & \textbf{Frequency} & \textbf{Spatial precision needed} \\ \hline
\textbf{Trend} & slow variation & low & rough \\ \hline
\textbf{Anomaly} & abrupt variation & high & fine \\ \hline
\end{tabular}}
\end{center}

Wavelets split image rows/blocks into approximation plus details:

\begin{itemize}
  \item approximation: low-resolution trend,
  \item detail: high-frequency anomalies, often sparse.
\end{itemize}

\subsection{1D Filter Bank}

Wavelet transform is implemented by analysis filters plus downsampling.

Analysis:

\[
x[k]\xrightarrow{h_0}\tilde{c}[k]\xrightarrow{\downarrow2}c[k]
\]

\[
x[k]\xrightarrow{h_1}\tilde{d}[k]\xrightarrow{\downarrow2}d[k]
\]

\begin{itemize}
  \item $h_0$: low-pass analysis filter, approximation coefficients.
  \item $h_1$: high-pass analysis filter, detail coefficients.
  \item Downsampling keeps one sample every two.
\end{itemize}

Synthesis:

\[
c[k]\xrightarrow{\uparrow2}\hat{c}[k]\xrightarrow{f_0}\bar{x}[k]
\]

\[
d[k]\xrightarrow{\uparrow2}\hat{d}[k]\xrightarrow{f_1}\bar{x}[k]
\]

Interpolation:

\[
\hat{c}[k]=
\begin{cases}
c[k/2] & k\ \text{even}\\
0 & k\ \text{odd}
\end{cases}
\]

\subsection{Perfect Reconstruction}

A useful wavelet filter bank should allow perfect reconstruction (possibly with delay) after analysis and synthesis.

In $z$ domain:

\[
\tilde{C}(z)=H_0(z)X(z)
\]

Decimation:

\[
C(z)=\frac{1}{2}
\left[
\tilde{C}(z^{1/2})+\tilde{C}(-z^{1/2})
\right]
\]

Interpolation:

\[
\hat{C}(z)=C(z^2)
\]

Output:

\[
\tilde{X}(z)
=
\frac{1}{2}[F_0(z)H_0(z)+F_1(z)H_1(z)]X(z)
+
\frac{1}{2}[F_0(z)H_0(-z)+F_1(z)H_1(-z)]X(-z)
\]

\begin{calloutbox}{red!70!black}{Perfect reconstruction conditions}
\textbf{No distortion:}

\[
T(z)=H_0(z)F_0(z)+H_1(z)F_1(z)=2z^{-\ell}
\]

\textbf{Aliasing cancellation:}

\[
A(z)=H_0(-z)F_0(z)+H_1(-z)F_1(z)=0
\]

Together:

\[
\tilde{X}(z)=z^{-\ell}X(z)
\]

so the output is the input delayed by $\ell$ samples.
\end{calloutbox}

\subsection{Orthogonality and Biorthogonality}

Orthogonal filter banks conserve energy:

\[
\sum_{k=-\infty}^{+\infty}x_k^2
=
\sum_{k=-\infty}^{+\infty}c_k^2
+
\sum_{k=-\infty}^{+\infty}d_k^2
\]

\textbf{Compression meaning:} reconstruction error equals quantization error in coefficient domain.

Biorthogonal filters are not strictly energy preserving, but they can be symmetric and close to orthogonal. In image compression, symmetry is often more important because it avoids boundary artifacts.

\subsection{Vanishing Moments}

\begin{calloutbox}{red!70!black}{Vanishing moments}
A wavelet/filter with $p$ vanishing moments has zero high-pass response to polynomials of degree $<p$.

It needs at least $2p$ taps.

Smooth image regions become nearly zero in detail subbands, increasing sparsity.
\end{calloutbox}

\subsection{Border Problem}

Wavelet theory assumes infinite signals, but images have finite support. Linear convolution with an $M$-tap filter expands an $N$-sample signal to $N+M-1$ samples.

Possible extensions:

\begin{center}\fitwidth{%
\begin{tabular}{|l|l|l|}
\hline
\textbf{Method} & \textbf{Benefit} & \textbf{Problem} \\ \hline
zero padding & simple & coefficient expansion and boundary artifacts \\ \hline
periodization & same number of coefficients & artificial boundary jumps \\ \hline
symmetrization & removes jumps & doubles period unless filters are symmetric \\ \hline
\end{tabular}}
\end{center}

\begin{calloutbox}{red!70!black}{Symmetry constraint}
The only FIR filter that is both orthogonal and symmetric, apart from trivial cases, is the Haar filter.

Image compression therefore prefers \textbf{biorthogonal symmetric filters} such as CDF 9/7 and CDF 5/3.
\end{calloutbox}

\subsection{Haar and CDF Filters}

Haar filters:

\[
h_0=[1\ 1],\qquad h_1=[1\ -1]
\]

\[
f_0=[1\ 1],\qquad f_1=[-1\ 1]
\]

Properties:

\begin{itemize}
  \item symmetric,
  \item orthogonal,
  \item one vanishing moment,
  \item only piecewise-constant approximation.
\end{itemize}

CDF filters:

\begin{center}\fitwidth{%
\begin{tabular}{|l|l|l|l|}
\hline
\textbf{Filter} & \textbf{Vanishing moments} & \textbf{Taps} & \textbf{Use} \\ \hline
\textbf{CDF 9/7} & 4 & 9/7 & JPEG2000 lossy, best R-D \\ \hline
\textbf{CDF 5/3} & 2 & 5/3 & JPEG2000 lossless, integer exact reconstruction \\ \hline
\end{tabular}}
\end{center}

If analysis has $p$ vanishing moments and synthesis has $\tilde{p}$, the filter needs at least:

\[
p+\tilde{p}-1
\]

taps.

\subsection{1D and 2D Multiresolution}

In 1D multiresolution, the filter bank is recursively applied only to the low-pass approximation $c_j$.

In 2D, apply 1D filters separably:

\begin{enumerate}
  \item filter rows,
  \item filter columns,
  \item downsample each direction.
\end{enumerate}

One level creates four subbands:

\begin{center}\fitwidth{%
\begin{tabular}{|l|l|l|}
\hline
\textbf{Subband} & \textbf{Filtering} & \textbf{Content} \\ \hline
\textbf{LL} & LP rows + LP columns & low-resolution approximation \\ \hline
\textbf{HL} & HP rows + LP columns & vertical details \\ \hline
\textbf{LH} & LP rows + HP columns & horizontal details \\ \hline
\textbf{HH} & HP rows + HP columns & diagonal details \\ \hline
\end{tabular}}
\end{center}

Recursive decomposition applies again to LL.

Optimal number of levels for images: \textbf{4 to 6}. More levels give diminishing returns because LL becomes too small and less well modeled as smooth trend.

\subsection{EZW}

\textbf{Embedded Zerotrees of Wavelet coefficients (EZW)} is a progressive wavelet coder:

\begin{itemize}
  \item sends largest coefficients first,
  \item uses bitplane thresholds,
  \item exploits inter-scale correlation,
  \item produces quality scalability,
  \item can be lossy-to-lossless with integer DWT and all bitplanes.
\end{itemize}

\begin{calloutbox}{red!70!black}{Zerotree principle}
If a wavelet coefficient is insignificant below threshold $T$, its descendants at finer scales in the same orientation are likely insignificant too.

A whole insignificant subtree can be coded with one \textbf{ZR} symbol.
\end{calloutbox}

EZW symbols:

\begin{center}\fitwidth{%
\begin{tabular}{|l|l|}
\hline
\textbf{Symbol} & \textbf{Meaning} \\ \hline
\textbf{SP} & significant positive, $c\geq T$ \\ \hline
\textbf{SN} & significant negative, $c\leq -T$ \\ \hline
\textbf{IZ} & isolated zero: insignificant but has significant descendant \\ \hline
\textbf{ZR} & zerotree root: insignificant and all descendants insignificant \\ \hline
\end{tabular}}
\end{center}

Thresholds:

\[
T_k=2^{n-k},
\qquad
n=\left\lfloor\log_2\max|c|\right\rfloor
\]

Each pass halves the threshold:

\[
T_{k+1}=\frac{T_k}{2}
\]

Dominant pass finds newly significant coefficients; refining pass sends the next bitplane of already significant coefficients.

\subsection{JPEG2000}

JPEG2000 addresses JPEG limitations:

\begin{center}\fitwidth{%
\begin{tabular}{|l|l|}
\hline
\textbf{JPEG issue} & \textbf{JPEG2000 response} \\ \hline
poor quality below 0.25 bpp & wavelet transform \\ \hline
blocking artifacts & global/tiled DWT instead of $8x8$ DCT \\ \hline
weak scalability & quality + resolution scalability \\ \hline
no random access & tiling and codeblock independence \\ \hline
imprecise rate control & EBCOT optimized truncation \\ \hline
no native lossy-to-lossless path & 5/3 integer DWT and scalable bitstream \\ \hline
\end{tabular}}
\end{center}

JPEG2000 structure:

\begin{itemize}
  \item \textbf{Tier 1:} DWT, fine quantization if lossy, arithmetic coding of codeblocks by bitplane.
  \item \textbf{Tier 2:} EBCOT organizes coded blocks into layers and truncates bitstreams optimally.
\end{itemize}

Lossy JPEG2000 uses CDF 9/7. Lossless JPEG2000 uses integer CDF 5/3.

\subsection{EBCOT}

\textbf{EBCOT} means \textit{Embedded Block Coding with Optimized Truncation}.

Workflow:

\begin{enumerate}
  \item split each DWT subband into codeblocks, e.g. $64x64$,
  \item independently arithmetic-code each codeblock by bitplanes,
  \item compute R-D curve for each codeblock,
  \item choose truncation points to satisfy target rate.
\end{enumerate}

Optimization:

\[
\min\sum_i D_i
\quad\text{s.t.}\quad
\sum_i R_i\leq R_{\text{tot}}
\]

Lagrangian:

\[
J=\sum_i(D_i+\lambda R_i)-\lambda R_{\text{tot}}
\]

\begin{calloutbox}{red!70!black}{EBCOT optimal truncation}
For each codeblock:

\[
\frac{\partial J_i}{\partial R_i}
=
\frac{\partial D_i}{\partial R_i}+\lambda=0
\]

At optimum:

\[
\frac{\partial D_i}{\partial R_i}=-\lambda
\]

All selected codeblock truncation points have the same R-D slope.
\end{calloutbox}

Multiple $\lambda$ values give multiple quality layers.

\subsection{Error Robustness}

Compressed streams are fragile because:

\begin{itemize}
  \item predictive coding propagates errors,
  \item variable-length coding loses synchronization,
  \item header errors can corrupt complete files.
\end{itemize}

JPEG2000 improves robustness through codeblock independence and bitstream organization.

\begin{center}\fitwidth{%
\begin{tabular}{|l|l|l|}
\hline
\textbf{Feature} & \textbf{JPEG} & \textbf{JPEG2000} \\ \hline
entropy stream & more global & independent codeblocks \\ \hline
error propagation & severe & localized \\ \hline
transform artifacts & blocking & no blocking \\ \hline
rate/resolution layers & limited & built in \\ \hline
\end{tabular}}
\end{center}

\section{Theory and Formulas}

\subsection{Wavelet Scaling and Fourier Interpretation}

Mother wavelet:

\[
\psi_{a,b}(t)
=
\frac{1}{\sqrt{a}}\psi\left(\frac{t-b}{a}\right)
\]

Fourier scaling:

\[
\mathcal{F}\{x(t/a)\}=|a|X(af)
\]

Thus large $a$ stretches the wavelet in time and compresses its frequency support. Small $a$ does the opposite.

\subsection{Perfect Reconstruction Matrix Form}

PR conditions can be written as:

\[
\begin{bmatrix}
H_0(z) & H_1(z)\\
H_0(-z) & H_1(-z)
\end{bmatrix}
\begin{bmatrix}
F_0(z)\\
F_1(z)
\end{bmatrix}
=
\begin{bmatrix}
2z^{-\ell}\\
0
\end{bmatrix}
\]

The modulation matrix must be invertible on the unit circle:

\[
\Delta(z)
=
H_0(z)H_1(-z)-H_1(z)H_0(-z)
\neq0,
\qquad |z|=1
\]

\subsection{EZW Zerotree Saving}

If a zerotree root starts at scale $n$ and has depth to $N$, one symbol can replace:

\[
1+4+4^2+\cdots+4^{N-n}
\]

individual insignificance symbols in 2D.

\subsection{JPEG2000 Exact Rate Control}

EBCOT can hit a target bitrate with very small error because each codeblock bitstream has many candidate truncation points:

\begin{itemize}
  \item early truncation: low rate, high distortion,
  \item later truncation: high rate, low distortion,
  \item layers collect truncation segments in increasing quality order.
\end{itemize}

\section{Visual Schemes}

\subsection{Time-Frequency Resolution}

\begin{center}\includegraphics[width=\linewidth]{ch05/time-frequency-short-window.png}\end{center}
Short windows localize time well but blur frequency.

\begin{center}\includegraphics[width=\linewidth]{ch05/time-frequency-long-window.png}\end{center}
Long windows localize frequency well but blur time.

\begin{center}\includegraphics[width=0.7\linewidth]{ch05/stft-rigid-tiling.png}\end{center}
STFT and block-DCT use fixed tiles, so all frequencies receive the same resolution tradeoff.

\begin{center}\includegraphics[width=0.7\linewidth]{ch05/wavelet-adaptive-tiling.png}\end{center}
Wavelets use narrow high-frequency tiles and wide low-frequency tiles.

\subsection{Trends and Anomalies}

\begin{center}\includegraphics[width=0.7\linewidth]{ch05/image-original-row.png}\end{center}
An image row contains both smooth trends and sharp localized changes.

\begin{center}\includegraphics[width=0.7\linewidth]{ch05/image-anomalies.png}\end{center}
High-pass wavelet details isolate anomalies such as edges and contours.

\begin{center}\includegraphics[width=\linewidth]{ch05/trend-detail-model.png}\end{center}
Approximation stores the trend; detail stores deviations from the trend.

\subsection{Filter Bank and Multiresolution}

\begin{center}\footnotesize\itshape Diagram (Mermaid source):\end{center}
\begin{verbatim}
flowchart LR
    X[Input x] --> LP[Low-pass h0]
    X --> HP[High-pass h1]
    LP --> D1[Downsample 2]
    HP --> D2[Downsample 2]
    D1 --> C[Approximation c]
    D2 --> D[Detail d]
\end{verbatim}

\begin{center}\includegraphics[width=\linewidth]{ch05/one-dimensional-mra.png}\end{center}
1D multiresolution recursively decomposes the low-pass approximation.

\begin{center}\includegraphics[width=\linewidth]{ch05/multilevel-decomposition.png}\end{center}
Analysis filter banks generate approximation and detail streams across scales.

\begin{center}\includegraphics[width=\linewidth]{ch05/multilevel-reconstruction.png}\end{center}
Synthesis filter banks invert the decomposition from coarse scale back to the original signal.

\subsection{2D DWT}

\begin{center}\includegraphics[width=\linewidth]{ch05/two-dimensional-dwt-subbands.png}\end{center}
One 2D DWT level produces LL, HL, LH, and HH subbands.

\begin{center}
\includegraphics[width=0.633\linewidth]{ch05/synthetic-square.png}\end{center}
Synthetic square image is simple but has sharp edges.

\begin{center}\includegraphics[width=\linewidth]{ch05/synthetic-square-dwt.png}\end{center}
DWT places smooth content in LL and edges/corners in detail subbands.

\subsection{EZW and JPEG2000 Schemes}

\begin{center}\footnotesize\itshape Diagram (Mermaid source):\end{center}
\begin{verbatim}
flowchart TD
    A[DWT coefficients] --> B[Initial threshold T0]
    B --> C[Dominant pass]
    C --> D{Coefficient class}
    D --> SP[SP / SN]
    D --> IZ[IZ]
    D --> ZR[ZR subtree]
    SP --> E[Significant list]
    E --> F[Refining pass]
    F --> G[Halve threshold]
    G --> C
\end{verbatim}

\begin{center}\footnotesize\itshape Diagram (Mermaid source):\end{center}
\begin{verbatim}
flowchart LR
    A[Image tile] --> B[DWT 9/7 or 5/3]
    B --> C[Subbands]
    C --> D[Codeblocks]
    D --> E[Bitplane arithmetic coding]
    E --> F[R-D curves]
    F --> G[EBCOT truncation]
    G --> H[Quality / resolution layers]
    H --> I[JPEG2000 bitstream]
\end{verbatim}

\begin{calloutbox}{blue!60!black}{Missing visual}
Source uses placeholders for EBCOT codeblocks, R-D truncation, JPEG/JPEG2000 comparison, and error-robustness figures. Mermaid schemes above cover the processing logic.
\end{calloutbox}

\section{Examples}

\begin{calloutbox}{green!55!black}{Trends vs anomalies}
\textbf{Setup:} an image row has smooth regions plus edges.

\textbf{Wavelet effect:} low-pass branch captures slowly varying intensity; high-pass branch activates mostly at edges.

\textbf{Takeaway:} smooth areas generate few detail coefficients, so they compress well.
\end{calloutbox}

\begin{calloutbox}{green!55!black}{Haar transform intuition}
For a pair $(x_0,x_1)$, Haar approximation and detail are proportional to:

\[
c=x_0+x_1,\qquad d=x_0-x_1
\]

If the pair is smooth, $x_0\approx x_1$, so $d\approx0$.

\textbf{Takeaway:} Haar is simple edge/detail extraction, but only has one vanishing moment.
\end{calloutbox}

\begin{calloutbox}{green!55!black}{2D DWT subbands}
\textbf{LL:} blurred low-resolution image.

\textbf{HL:} responds to vertical details.

\textbf{LH:} responds to horizontal details.

\textbf{HH:} responds to diagonal details.

\textbf{Takeaway:} DWT separates image geometry by scale and orientation.
\end{calloutbox}

\begin{calloutbox}{green!55!black}{EZW encoding on a 4x4 coefficient matrix}
Coefficients:

\[
\begin{bmatrix}
26 & 6 & -13 & 10\\
-7 & 3 & 6 & 4\\
4 & -3 & 3 & -3\\
2 & -2 & -2 & 0
\end{bmatrix}
\]

Initial threshold:

\[
T=16
\]

Since $26\geq16$, first coefficient is \textbf{SP}. Coefficients below threshold whose descendants are also below threshold become \textbf{ZR}.

First dominant pass:

\begin{verbatim}
SP ZR ZR ZR
\end{verbatim}

Refining sends next bit of already significant coefficients.

\textbf{Takeaway:} EZW sends coarse significance first, then progressively refines.
\end{calloutbox}

\begin{calloutbox}{green!55!black}{EZW decoding estimate}
If a coefficient becomes significant positive at threshold $T=16$, decoder knows:

\[
16\leq c<32
\]

Initial estimate is midpoint:

\[
\hat{c}=24
\]

Refining bits shrink the uncertainty interval by factor 2 at each pass.
\end{calloutbox}

\begin{calloutbox}{green!55!black}{JPEG2000 vs JPEG at low rate}
\begin{center}\fitwidth{%
\begin{tabular}{|l|l|l|}
\hline
\textbf{Rate} & \textbf{JPEG} & \textbf{JPEG2000} \\ \hline
1.0 bpp & good & good \\ \hline
0.5 bpp & visible blocking & clean \\ \hline
0.3 bpp & strong blocking & slight blur \\ \hline
0.2 bpp & very poor & acceptable \\ \hline
0.1 bpp & often unusable & usable \\ \hline
\end{tabular}}
\end{center}

\textbf{Takeaway:} wavelet coding degrades by blur/ringing rather than block artifacts.
\end{calloutbox}

\begin{calloutbox}{green!55!black}{Error robustness}
At bit error probability $p_E=10^{-3}$, JPEG can suffer catastrophic visual degradation because entropy decoding can desynchronize and errors propagate.

JPEG2000 confines much damage to independent codeblocks, so degradation is more local and graceful.
\end{calloutbox}

\begin{calloutbox}{green!55!black}{Method comparison}
\begin{center}\fitwidth{%
\begin{tabular}{|l|l|l|l|l|l|}
\hline
\textbf{Technique} & \textbf{Transform} & \textbf{Coding} & \textbf{Scalability} & \textbf{Lossless path} & \textbf{Blocking} \\ \hline
JPEG & $8x8$ DCT & Huffman/arithmetic & no & no & yes \\ \hline
EZW & DWT & symbols + arithmetic & quality & yes with integer DWT & no \\ \hline
JPEG2000 & DWT 9/7 or 5/3 & EBCOT arithmetic & quality, resolution, ROI & yes with 5/3 & no \\ \hline
\end{tabular}}
\end{center}
\end{calloutbox}
